\subsection{Coupled vacuum families from two through twelve loops}
\label{sec:higher-loop-generation}
To supply examples at every loop order considered in the model study,
we added 100 mass profiles at each order from two through twelve.
These profiles define 1,100 additional family instances. They use
coupled vacuum integrals whose master expansions follow from mass
partial fractions and IBP identities, which allows us to retain exact
reference answers as the number of loops increases. The profiles vary
the masses within a restricted set of graphs, so the family count
measures mass assignments as well as graph structure.

For this extension we use Euclidean loop momenta and define
\begin{equation}
 V^{(L)}_{\boldsymbol b}=
 \int\prod_{i=1}^{L}\frac{d^d\ell_i}{\pi^{d/2}}\,
 \frac{1}{\ell_1^2\cdots\ell_L^2}
 \prod_{j=1}^{R}\frac{1}{(Q^2+\mu_j)^{b_j}},
 \qquad
 Q=\sum_{i=1}^{L}\ell_i,\qquad \mu_j=\rho_j\mu .
 \label{eq:coupled-vacuum}
\end{equation}
Here $L$ is the number of loops, $R$ is the number of distinct massive
denominators, $\mu=m^2>0$ is a common squared-mass scale and
$\rho_j$ is a positive mass ratio. The nonnegative integer $b_j$
is the power of the $j$th massive denominator, and
$B=\sum_{j=1}^{R}b_j\geq1$ is their total power. Every massless
denominator has power one in the generated domain. The full integral
index vector is $(1,\ldots,1,b_1,\ldots,b_R)$, with $L+R$ entries.
Equation~\eqref{eq:coupled-vacuum} specifies the normalization and signs
used for this Euclidean subcollection; the earlier families retain
their stated Minkowski convention.

The common momentum $Q$ couples the loop integrations. A graph
representation has $L$ parallel massless edges between two vertices
and a chain of $R$ massive edges joining the same vertices. It has
$L+R$ edges and $R+1$ vertices, giving $L$ independent loops.
All massive denominators depend on $Q^2$, which gives algebraic
relations between them. This definition is a parent family for partial
fractions with fixed massless powers; it does not specify a complete
independent denominator and numerator basis for a generic
$L$-loop scattering integral.

For profile number $p=1,\ldots,100$, we choose
\begin{equation}
 R=2+((p-1)\bmod3),\qquad
 \rho_1=1,\qquad
 \rho_{j+2}=2+3p+j(p+2),\quad j=0,\ldots,R-2 .
 \label{eq:higher-loop-masses}
\end{equation}
The ratios are distinct within each profile, and $\rho_1=1$ fixes a
common rescaling. This gives 34 profiles with two massive
denominators, 33 with three and 33 with four at each loop order.
A twelve-loop family with four massive denominators has sixteen
indices, which fits the existing model representation.

\subsection{Exact masters and source equations for the added families}
The dependence on the common momentum makes the master basis
explicit. Write $\boldsymbol e_j$ for the unit vector in massive
index position $j$ and choose
\begin{equation}
 M_j^{(L)}=V^{(L)}_{\boldsymbol e_j},
 \qquad j=1,\ldots,R .
 \label{eq:higher-loop-masters}
\end{equation}
To see why these integrals span the generated domain, first use
$D_i-D_j=\mu_i-\mu_j$, where $D_j=Q^2+\mu_j$. If $b_i,b_j>0$,
this identity gives
\begin{equation}
 (\mu_i-\mu_j)V^{(L)}_{\boldsymbol b}
 -V^{(L)}_{\boldsymbol b-\boldsymbol e_i}
 +V^{(L)}_{\boldsymbol b-\boldsymbol e_j}=0 .
 \label{eq:higher-loop-partial-fraction}
\end{equation}
Repeated use removes products of distinct massive denominators.
This equation is an algebraic mass relation and carries no IBP seed.

The remaining powers follow from the sum of the $L$ IBP operators
$\partial_{\ell_i^\nu}\ell_i^\nu$. Applying this sum to the integrand
of Eq.~\eqref{eq:coupled-vacuum} gives
\begin{equation}
 0=[Ld-2(L+B)]V^{(L)}_{\boldsymbol b}
   +2\sum_{j=1}^{R}b_j\mu_j
                 V^{(L)}_{\boldsymbol b+\boldsymbol e_j}.
 \label{eq:higher-loop-euler}
\end{equation}
Each such equation retains its actual seed
$(1,\ldots,1,\boldsymbol b)$ and the contributing derivative operators.
For one massive denominator this becomes
\begin{equation}
 V^{(L)}_{k\boldsymbol e_j}
 =\left[\prod_{h=1}^{k-1}
       \frac{L+h-Ld/2}{h\mu_j}\right]M_j^{(L)},
 \qquad k\geq1 .
 \label{eq:higher-loop-dots}
\end{equation}
The empty product at $k=1$ equals one. Equations
\eqref{eq:higher-loop-partial-fraction} and
\eqref{eq:higher-loop-dots} therefore reduce every generated integral
to the stated masters with coefficients rational in $d$ and $\mu$.

An analytic expression also checks the basis. Set $D=d/2$,
$\alpha=L(D-1)$ and $c=L-(L-1)D$. Repeated integration of the
massless subgraph at fixed $Q$ gives
\begin{equation}
 C_L(D)(Q^2)^{-c},\qquad
 C_L(D)=\frac{\Gamma(c)\Gamma(D-1)^L}{\Gamma(\alpha)} .
 \label{eq:higher-loop-convolution}
\end{equation}
The remaining radial integration uses
$\int d^dQ/\pi^D=\Gamma(D)^{-1}\int_0^\infty
dx\,x^{D-1}$ for a rotationally invariant integrand, where $x=Q^2$.
It gives
\begin{equation}
 M_j^{(L)}=
 \mu_j^{\alpha-1}
 \frac{\Gamma(c)\Gamma(D-1)^L\Gamma(1-\alpha)}{\Gamma(D)} .
 \label{eq:higher-loop-master-value}
\end{equation}
Here $\Gamma$ is the Euler gamma function. These manipulations use
the standard analytic continuation of dimensionally regularized
integrals~\cite{Weinzierl2022}. At generic $d$, distinct $\rho_j$
give distinct exponential functions of $d$ in
$\rho_j^{\alpha-1}$. Their linear independence over rational
functions of $d$ establishes independence of the $R$ masters in this
restricted domain. After the common factors cancel, a proposed
relation would combine exponentials with distinct growth rates.
At fixed generic $\mu$, rational coefficients grow or decay only
as powers of $d$, so they cannot cancel the largest exponential
unless its coefficient vanishes. Repeating this argument removes
each coefficient. The extension contains 299 master entries per
loop order and 3,289 added entries in total.

We also obtain the master coefficients without following the selected
row operations. With $z=Q^2/\mu$, the rational identity
\begin{equation}
 \prod_{i=1}^{R}(z+\rho_i)^{-b_i}
 =\sum_{j:b_j>0}\sum_{k=1}^{b_j}
                 A_{jk}(z+\rho_j)^{-k}
 \label{eq:higher-loop-oracle}
\end{equation}
defines rational numbers $A_{jk}$. Expanding the other factors
around $z=-\rho_j$ gives them directly:
\begin{equation}
 A_{jk}=[w^{\,b_j-k}]
           \prod_{i\ne j}(\rho_i-\rho_j+w)^{-b_i}.
 \label{eq:higher-loop-taylor}
\end{equation}
Here $w=z+\rho_j$ is the expansion variable, and $[w^n]$ means
the coefficient of $w^n$ in this Taylor
series. Combining these coefficients with
Eq.~\eqref{eq:higher-loop-dots} yields the independent reference
\begin{equation}
 V^{(L)}_{\boldsymbol b}
 =\mu^{1-B}\sum_{j:b_j>0}\sum_{k=1}^{b_j}A_{jk}
   \left[\prod_{h=1}^{k-1}
             \frac{L+h-Ld/2}{h\rho_j}\right]M_j^{(L)} .
 \label{eq:higher-loop-reference}
\end{equation}
The factor $\mu^{1-B}$ ensures that both sides have mass dimension
$Ld-2L-2B$. This reference allows us to check the reduction path
through a separate calculation.

\subsection{Reverse target construction and the expanded data division}
We supply every Euler seed with $1\leq B\leq5$ and retain the
summed Euler identity in Eq.~\eqref{eq:higher-loop-euler} as its
source equation together with all mass partial-fraction identities
through $B=6$. Scalar target roots have $2\leq B\leq5$.
The source pools contain 35, 160 or 545 equations for $R=2,3,4$,
respectively. These counts stay fixed as $L$ changes because the
massless powers and massive-index bounds stay fixed. The additional
loop orders therefore test this controlled family structure, while
large generic IBP systems require wider denominator and numerator
domains.
After obtaining a checked scalar relation, we reverse its steps to
construct the corresponding integral from its master combination.
For expressions containing several roots, we first combine their
known master expansions and then apply the recorded inverse steps.
Each inverse step adds a stated multiple of an exact zero identity,
so the master coefficients remain known throughout the construction.

For roots with total powers $B_\nu$, we use target weights of the form
\begin{equation}
 {\cal T}=\sum_\nu
   \kappa_\nu d^{\eta_\nu}\mu^{B_\nu-B_*}
                       V^{(L)}_{\boldsymbol b_\nu},
 \qquad
 B_*=\max_\nu B_\nu,\qquad \eta_\nu\in\{0,1\}.
 \label{eq:higher-loop-target-weights}
\end{equation}
Here $\kappa_\nu$ is a nonzero rational number. The powers of $\mu$
give every term the same mass dimension. Expressions use up to three
distinct roots and include a root with every massive denominator
present and at least one additional denominator power. Individual
targets also sample lower sectors. We retain 128 distinct targets
per family, giving 140,800 additional targets.

For each target we retain its master coefficients, exact reverse
construction, source equations, sufficient IBP seed set and a
forward elimination trace. The trace records all allowed source
rows, the chosen pivot, the coefficients before and after the
operation, and the arithmetic operation counts. Mass
partial-fraction steps remain separately identified, so a seed label
always refers to an actual IBP seed. The reference path minimizes
the next symbolic row's number of nonmaster terms, then its total
support and immediate operation count, with a fixed final tie rule.
This supplies a reproducible valid path without establishing a
globally minimal sequence.

We verify every source equation against
Eq.~\eqref{eq:higher-loop-reference} and separately check the
partial-fraction formula and the Euler derivative before integration. We replay the symbolic constructions
and forward reductions, then independently replay all saved records
at three finite-field samples. A candidate target is rejected if a
reference pivot vanishes at one of those samples. The checks also
verify connected graphs and their loop counts. At $d=2+1/L$ the
radial master integral converges; numerical quadrature then checks
Eq.~\eqref{eq:higher-loop-master-value} at both ends of the loop-order
range.

The complete collection now contains 2,100 family instances,
12,967 master entries and 381,272 targets. It contains 997 one-loop
families, 103 two-loop families and 100 families at each order from
three through twelve. We retain the original grouped model
assignment and assign 70 added mass profiles to training and 30 to
testing. A profile has the same assignment at every loop order,
giving 266,890 training targets in 1,463 families and 114,382
test targets in 637 families. Shared lower-sector integrals can occur
across different profiles. These tests measure transfer between the
specified mass profiles with every loop order represented during
training; they do not establish transfer to unseen loop orders or
general high-loop topologies. The accompanying records contain the
complete definitions and checks~\cite{HigherLoopIntegralData}.
